\documentclass[12pt]{article}

\usepackage{amsmath}
\usepackage[style=authoryear,backend=biber]{biblatex}
\addbibresource{infillPuzzle}


\begin{document}

\title{Where does the infill housing go?\\Addition versus division}
\author{Thomas Davidoff, Robert Helsley, and Tsur Somerville\thanks{thomas.davidoff@sauder.ubc.ca; robert.helsley@sauder.ubc.ca; tsur.somerville@sauder.ubc.ca}\\Sauder School of Business, University of British Columbia\\\textbf{Under continual revision.} Updated version is available at:\\ https://github.com/tomdavidoff/infillPuzzle/blob/main/text/main.pdf}
  \maketitle

\onehalfspacing
\begin{abstract}
	
Jurisdictions across Canada and the U.S. have recently allowed ``missing-middle'' infill housing, such as accessory dwelling units and multiplexes, in residential zones formerly restricted to single-family use. Existing studies document that infill housing is less prevalent in more expensive neighbourhoods within metro areas. This is puzzling because the right to add space should be more valuable in neighbourhoods where land prices are higher. To resolve this contradiction, we develop a parsimonious model that distinguishes between two components of missing-middle reform: ``addition'' (adding to allowable square footage) and ``division'' (permitting the subdivision of a fixed amount of square footage into multiple units).  When density is not allowed to vary much with location, more affluent households will live in more desirable locations, so that division may be less valuable in better locations. Recent zoning reforms in Vancouver illustrate these mechanics: when duplexes became permitted in formerly single-owner zones, they were not permitted any additional density relative to single family homes and were strongly sorted into lower-priced locations. A subsequent reform downgraded single-owner homes' allowable density and was associated with relatively positive sorting of duplexes and newly allowed multiplexes into pricier locations.
par\bigskip
\textbf{Keywords:} housing supply and markets, government policy, regulation, state and local governmenjurisdictions
\end{abstract}

\section{\label{sec:intro}Introduction}

As concerns over housing affordability have grown in political salience over the last decade, a growing number of U.S. and Canadian jurisdictions have introduced reforms to allow ``missing middle'', ``infill'', or ``gentle density'' housing types in formerly single-family-only residential zones. This housing might take the form of additional density that does not dramatically change the urban fabric such as accessory dwelling units (Vancouver's laneway homes) or ``multiplex'' structures with a range of two to six homes on single lot, or more noticeable changes in urban form in the missing-middle's densest form such as row and townhouses or three- to four-storey apartment buildings. As has been emphasized, e.g. by \textcite{bruecknerSingh} and \textcite{HsiehMoretti}, restrictive zoning causes deadweight loss most where real estate is most valuable. It is thus natural to expect that the option to densify urban land would be exercised most commonly in the neighborhoods where land and structure prices per unit are highest.

The puzzle with missing middle uptake is that recent empirical studies have documented that infill housing is, in fact, less prevalent in the most expensive neighbourhoods in jurisdictions that have upzoned. We present a parsimonious model of density choice that shows how where redevelopment occurs can depend on the extent to which policy changes allow the addition of greater density on a lot (``addition'') and the allowance for more units within that density (``division''). These two effects likely work in opposite directions with respect to the relationship between infill prevalence and neighbourhood price levels. In the model, the extent that infill policies allow more square feet to be built on the same lot (addition) lead to greater take-up where the price per square foot of structure net of construction costs is greatest. For a given allowed density, allowing more (i.e. smaller) units to be created out of the same square footage (division) will result in negatively sorted take-up if wealthier households with a greater elasticity of willingness to pay for square footage sort into more expensive single-family neighbourhoods. That positive sorting is theoretically likely: in standard monocentric models affluent households may suburbanize due to flat bid-rent curves generated by income-elastic demand for housing,\footnote{See \textcite{Wheaton77} and earlier work of Muth, Alonso, Mills, and Beckmann cited therein.} but with single-family zoning leading many households to a fixed (over-) consumption of lot size, the normality of demand for location quality should generate positive sorting.

This paper studies patterns of redevelopment from two very modest, but broad upzonings in Vancouver, BC, in Canada. These upzonings represent different mixes of addition and division. In the first case, allowing duplexes (division) without an increase in allowable floor area (addition). The second policy change allowing up to 8 units on single lot included addition.\footnote{We focus on the reforms that allowed for additional ownership units rather than those that allowed for accessory dwelling units, which are rental only.} We evaluate the spatial correlation between neighbourhood price levels and the prevalence of infill housing permits following the different upzoning. The first policy allowing duplexes without additional floor area has been taken up predominantly in less expensive East Side neighbourhoods. Relatively recently-allowed multiplexes, that allow a bit more density have also been concentrated on the East Side; a quantitative exercise does not make a clear prediction of the siting given the offsetting addition and division effects. (Our empirical evidence here is on the first reform, analysis of the multiplex reform and similar reforms in Portland and Minneapolis are work in progress).

The remainder of the paper is structured as follows. In Section \ref{sec:lit} we review the literature on upzoning. Section \ref{sec:model} presents our model. We follow this with a description of the institutional and policy setting in Section \ref{sec:policies}. Our empirical analysis is in Section \ref{sec:results}. Section \ref{sec:conclusion} concludes. The theoretical model consider small changes to a single type of housing in an open city with two income groups. With fixed home sizes, under minimal assumptions, households are sorted by income, with higher income households in better locations. An increase in allowable space, holding the number of units constant, makes the high-income area expand because of the high-income households' greater willingness to pay for space. An increase in the number of units allowed, holding the total space constant, makes the low-income area expand, again because of the high-income households' greater willingness to pay for space. The model thus predicts by analogy that reforms allowing more space per unit will be taken up predominantly in high-income neighbourhoods, while reforms allowing more units per lot holding space constant will be taken up predominantly in low-income neighbourhoods. The spatial sorting of redevelopment options that mix addition and division (such as Vancouver's multiplex reform of 2023) is ambiguous. The model predicts ``by analogy'' because we cannot easily predict whether sorting will ``flip'' in the presence of a division reform: it is possible that a large reduction in space per unit would trigger an equilibrium in which poorer households live in smaller units closer to downtown, as is possible in the standard monocentric model. 

\section{\label{sec:lit}Related Literature}

This paper draws from several different threads of urban economics. The first is an older literature on urban growth under a monocentric model. The second is the literature on residential sorting, and the third is more current work on the effects of land use regulation reforms. Our contribution is to  examine the sensitivity of residential redevelopment to the form of zoning in the presence of households sorting across space. 

Redevelopment more broadly has been long studied in the literature, beginning with the dynamic models of urban growth such as \textcite{Brueckner_1981}, \textcite{HochmanPines_1980}, and \textcite{Wheaton_1982}. With durable but replaceable capital in a monocentric model and no other neighborhood specific amenities, \textcite{Wheaton_1982} finds that redevelopment occurs with substantial increases in the capital-to-land ratio (density and or structure quality) and that, under standard conditions occurs primarily in the urban core. Those papers have no redevelopment frictions caused by local land use policies. Introducing the restrictions as frictions is important as they have been identified by economists as the primary cause of limits on housing supply that increase housing prices (see \textcite{GyourkoMolloy2015} and \textcite{Molloy2020} for comprehensive reviews of this literature).

Over the last decade there have been broad rezonings enacted in cities around the world. The empirical study of the effects of these large scale upzonings on new supply finds support for the positive effect of relaxing density restrictions on redevelopment at higher densities. 

\textcite{Anagol_etal2021} analyze a broad-based 2016 zoning reform in S\~{a}o Paulo, Brazil, that allowed an average increase in the maximum allowed built area ratio (BAR) - analogous to floor area ratio (FAR) or floor space ratio (FSR) - from 1.54 to 2.09. Over six years they note a 66 percent increase in multi-family permit applications for treated blocks compared to the controls (or .012 permits per block per year) resulting in an aggregate 1.9 percent increase in the stock of housing and a 0.5 percent decline in prices.

Like S\~{a}o Paulo, Auckland, NZ also introduced a single broad upzoning. Implemented by the Auckland Council (the government body administering the metro area) in late 2016, the Auckland Unitary Plan (AUP), increased the allowed number of units and density on about 75 percent of the metro area's residential land (about 90 percent of the upzoned area allowed for three units per parcel, with the remaining 10 percent at higher densities). \textcite{GreenawayPhillips2023} identify a large increase in new construction (equal to an estimated 4.1 percent of the existing housing stock over five years) attributable to the policy changes in the AUP. 

Both \textcite{Anagol_etal2021} and \textcite{GreenawayPhillips2023} study relatively short-run effects, within a five-year period following the rezoning. \textcite{BuchlerLutz2024} are able to examine longer-term effects studying Swiss zoning reforms that increased densities over a multi-decade period starting in the 1990s. They study rolling increases in the allowed maximum density (with a median 30\% increase in density) across the Canton of Zurich from 1995 to 2020 using an event study approach for each hectare in Zurich. For a five-ten year post-reform period they find larger effects than in \textcite{Anagol_etal2021}: a 9 percent increase in living space and housing units in up-zoned compared to not up-zoned areas. As well, uptake is greater where existing buildings are already at the maximum allowed density, i.e. zoning is binding. 

While not the entire city, New York, NY introduced a series of neighbourhood-level rezonings affecting approximately 40 percent of the city between 2004 and 2013 (\textcite{Liao2023}). Using a border discontinuity design for local increases in FAR, \textcite{Liao2023} finds an overall 4 percent increase in housing supply over seven years, and 8 percent for areas that were treated with the largest increases in FAR. Using the same zoning reform, \textcite{Peng2023} creates a spatial model with worker residential mobility and finds as a counterfactual to the observed outcomes a lower effect of a 0.7 percent long-run increase in the stock. 

\textcite{Rollet_2025} follows more explicitly in the spirit of models of redevelopment in including patterns of existing development in the assessment of uptake in the upzoning reforms. He creates a dynamic model of the redevelopment choices made by developers within a dynamic spatial equilibrium for households and firms. He finds that the large fixed costs of redevelopment, especially in dense areas, limits the short-run effects of rezoning. In simulations, after 40 years only 18\% of the allowed additional floor space would actually be added.  

In Minneapolis and Portland that enacted similar zoning changes to Vancouver, with the same greater emphasis on division than addition, studies find mixed effects on supply. \textcite{Dong2024a} assesses the Portland, OR upzoning. He finds that as a result of these changes, ``moderate middle'' housing increased from 13.4 percent to 44.7 percent of new building permits in single family zones, but does not provide an estimate of the total change in housing.\textcite{GuMunro_2025} use synthetic controls to test the effects of the Minnesota reforms, finding lower prices and rents, but surprisingly no change in permits.

Our work looks at how sorting can affect rezoning uptake. Starting with \textcite{Tiebout1956} and with the empirical foundation for effects on house prices in \textcite{Rosen1974}, numerous papers have examined sorting by households across space. Sorting on the basis of amenities has been general (\textcite{Bayer_etal2004}), but often in a US context focused on the sorting effects of preferences for school quality. More recent work, as in \textcite{Couture_etal2024} and \textcite{AlmagoDomingues2025}, has extended this to endogenous amenities, such as gentrification. Our model, like that in \textcite{BayerMcMillan2012} focuses on the joint preference for amenities and house type, though in a more focused framing that identifies the effects of the levers of zoning on sorting. 

Several authors have identified geographic patterns of redevelopment activity following zoning reforms. \textcite{MarantzElmendorfKim2023_ADUReforms} find that accessory dwelling unit (ADU) placement in California is widespread, but concentrated in convenient, low- and moderate-rent neighbourhoods, and not in more expensive neighbourhoods. Similar findings are in \textcite{BruecknerThomaz2024}, who find that ADU permitting rates are lower in higher-priced neighbourhoods, even after controlling for income and other factors. Huang et al (2023) analyze the construction of duplexes and triplexes in Seattle following a 2019 upzoning reform and find that these housing types are more likely to be built in lower-priced neighbourhoods. Finally, \textcite{DongHansz_2019} examine the impact of Portland's Residential Infill Project (RIP) on housing construction patterns and find that the new zoning has led to more infill development in lower-priced areas. 

While the work cited above examines upzoning in individual cities, \textcite{Stacy_etal2023} create a cross-city panel to study the effects of relaxing zoning restrictions on new construction more broadly. They conduct newspaper text searches for approximately 180 land use regulatory changes between 2000 and 2019 in eight US metro areas. About half (98) of these lessen restrictions through either allowing ADUs (35) or increasing FARs (15). If they observe a relaxation in any measure, they observe the discrete characteristic of having loosened restrictions is associated with a 0.8 percent increase in residential postal addresses in the jurisdictions that loosen restrictions, when compared to jurisdictions that do not. In terms of geographic patterns, they find the effects come mainly from above-median rent areas.

We study zoning that allows some combination of increased units or allowed floor area on a given lot. With undeveloped lots, an alternative to allowing duplexes would be halving lot sizes. Using data from Atlanta to calibrate a model, \textcite{Ma_2025} estimates the effect of this type of forced upzoning. In her data, the binding constraint is not allowed density, but division, where large lot zoning yields large single family homes in well off areas. In her model she allows for demand heterogeneity in preferences for unit type and locations. Regulatory costs cause developers to build large units on large lots rather than trying to get permission to subdivide and built multiple smaller units. She finds that forcing upzoning results in increased supply as when a single large unit cannot be built, more smaller units are constructed instead. Similarly, \textcite{vonBergmannLauster} find large empirical gains in rare cases of lot-splitting in Vancouver.


\section{\label{sec:model} A model of sorting and density regulations}

Locations in a city are vertically differentiated by quality $\theta$. $\theta$
can be determined by multiple characteristics, but there is agreement on the
ranking. At a location with quality $\theta$, price per square foot of a home
in an Nplex (a collection of $N$ homes), each of which has square footage
$s(N)$ is $p(N,\theta)s(N)$. We assume that builders are always at a corner
solution of building up to $s(N)$ square feet per unit. The function
$p(N,\theta)$ includes the relationship between size and number of units
implicitly in $N$. The cost per square foot of an Nplex is $c(N)$, with $c'\geq
0$ because of relatively fixed plumbing and utility costs and fire separation.

The profit from building an Nplex at location $\theta$ is thus:

\begin{equation}
	\pi(N,\theta) = \left[p(N,\theta)-c(N)\right]s(N)
\end{equation}

Defining $\Delta(n,\theta)$ as the difference in profit from developing $n+1$ instead of $n$ units, $\Delta(n,\theta) \equiv \pi(n+1,\theta)-\pi(n,\theta)$, we can identify the effect of location on that gain in profits as an effect on dividing the same amount of space into more units and the effect on location on the benefit to adding more square footage:

\begin{equation}
	\label{eq:dDelta}
	\begin{split}
		\frac{\partial \Delta(n,\theta)}{\partial \theta} = \underbrace{\frac{\partial\left[p\left(n+1,\theta\right)-p\left(n,\theta\right)\right]}{\partial \theta}s(n)}_{\text{division effect}} \\ + \underbrace{\frac{\partial p\left(n+1,\theta\right)}{\partial \theta}\left[s(n+1)-s(n)\right]}_{\text{addition effect}}.
	\end{split}
\end{equation}

Assuming positive sorting by income into neighbourhoods where $n$ or $n+1$ units are allowed, the first term may be negative. When location is better ($\theta$ greater), higher income households' elasticity of willingness to pay for homes with respect to space will be greater, so the percentage difference in price $p(n+1,\theta)/p(n,\theta)$ will be smaller. This does not unambiguously sign the division effect as negative, however, because the absolute difference could be larger given a higher base willingness to pay for better locations. Also, given higher income households' greater demand for larger and for single family homes, there can be ranges of $\theta$ over which there is negative income sorting: lower income households living in large Nplexes may outbid high income households who prefer single family homes.



\begin{equation}
	\frac{\partial \pi(N,\theta)}{\partial N} = \left[p(N,\theta)-c(N)\right]s'(N) + \left[\frac{\partial p(N,\theta)}{\partial N}-c'(N)\right]s(N).
\end{equation}

Differentiating with respect to $\theta$:

\begin{equation}
	\frac{\partial^{2} \pi(N,\theta)}{\partial N \partial \theta} = \underbrace{\frac{\partial^{2} p(N,\theta)}{\partial N \partial \theta}s(N)}_{\text{division effect}} + \underbrace{\frac{\partial p(N,\theta)}{\partial \theta}s'(N)}_{\text{addition effect}}
\end{equation}

The sign of the first, the ``division effect'' term depends on income sorting at home size $N$. If higher income households live in better locations at $N$ units per lot, then we would expect $\frac{\partial^{2}\pi}{\partial N\partial \theta}<0$, as sharing space and smaller spaces are inferior goods. The second term, the ``addition effect'' is positive if $s'(N)>0$, that is, if choosing a larger number of units allows more square feet to be built. 

\section{The division effect in pre-reform Vancouver}
Prior to the 2018 Vancouver infill reform, duplex sales were concentrated \dots



\end{document}
